\documentclass[../../../include/open-logic-section]{subfiles}

\begin{document}

\olfileid{sth}{replacement}{ref}
\olsection{प्रतिस्थापन आणि प्रतिबिंबन}

\stagesinex{} च्या साहाय्याने प्रतिस्थापनाचे समर्थन करण्याचा आपला शेवटचा प्रयत्न एका सखोल आणि सुंदर निष्कर्षापासून सुरू होतो:\footnote{आठवण: स्पष्टपणे अन्यथा म्हटले नसेल, तर सर्व सूत्रांमध्ये प्राचले असू शकतात.}
%यात $x_1, \ldots, x_n$ यांसारखे वरची रेघ असलेले लेखन वापरून
%``$x_1, \ldots, x_n$'' संक्षिप्त लिहिले आहे:
\begin{thm}[प्रतिबिंबन रूपबंध]\ollabel{reflectionschema}
कोणत्याही सूत्र $\phi$ साठी:
\[
\forall \alpha \exists \beta > \alpha (\forall x_1 \ldots, x_n \in
V_\beta)(\phi(x_1, \ldots, x_n) \liff \phi^{V_\beta}(x_1, \ldots, x_n))
\]
\end{thm}
\noindent
\olref[sth][replacement][strength]{formularelativization} प्रमाणे, $\phi$ मधील प्रत्येक संख्यक $V_\beta$ या संचापुरता मर्यादित केल्यावर $\phi^{V_\beta}$ मिळते. म्हणून सहज अर्थाने प्रतिबिंबन असे सांगते: $\phi$ संपूर्ण पदानुक्रमात सत्य असेल, तर पदानुक्रमाच्या हवे तितक्या अनेक \emph{आरंभीच्या खंडांमध्ये}ही $\phi$ सत्य असते.

\citet{Montague1961} आणि \citet{Levy1960} यांनी दाखवले की (योग्यरीत्या मांडलेली) प्रतिस्थापन आणि प्रतिबिंबन ही $\Z$ च्या साहाय्याने समतुल्य आहेत; म्हणजे त्यांपैकी कोणतेही एक जोडल्यास $\ZF$ मिळते. (हे निष्कर्ष आपण \olref[sth][replacement][refproofs]{sec} मध्ये सिद्ध करू.) ही समतुल्यता लक्षात घेऊन, \stagesinex{} वरून प्रतिबिंबन आणि प्रतिस्थापनाचे समर्थन करण्याचा पुढील प्रयत्न करता येईल: \stagesinex{} असेल, तर पदानुक्रम खूपच उंच असायला हवा; इतका उंच की त्याविषयी आपण म्हणू शकतो अशा कोणत्याही गोष्टीने त्याच्या उंचीला बंध घालता कामा नये. हा विचार आपण असा समजू शकतो: कोणतेही वाक्य~$\phi$ संपूर्ण पदानुक्रमात सत्य असेल, तर ते पदानुक्रमाच्या हवे तितक्या अनेक आरंभीच्या खंडांतही सत्य असते. हेच प्रतिबिंबन आहे.

पुन्हा, प्रतिस्थापनाचे आंतरिक समर्थन देण्याचा हा खरोखर आशादायक प्रयत्न वाटतो. पण याविषयी सांगण्यासारखे येथे मावेल त्याहून खूप अधिक आहे. तो सफल होतो की नाही, हे आता तुम्ही स्वतः ठरवावे.\footnote{पुढे वाचायचे असल्यास \citet[95--100]{Incurvati2020} उपयुक्त ठरेल.}

\end{document}
